% BEGIN GENERATED ARTIFACT: prf:hypothetical-syllogism-validity-propositional-logic
\newpage
\phantomsection
\label{prf:hypothetical-syllogism-validity-propositional-logic}
\LRAProofFor{thm:hypothetical-syllogism-validity-propositional-logic}

\begin{remark*}[Return]
\hyperref[thm:hypothetical-syllogism-validity-propositional-logic]{Return to Theorem (Hypothetical Syllogism is Valid).}
\end{remark*}

\begin{theorem*}[Hypothetical Syllogism is Valid]
The argument form
\[
\varphi\to\psi,\qquad \psi\to\chi \quad\therefore\quad \varphi\to\chi
\]
is valid.
\end{theorem*}

\begin{proof}
\textbf{Professional Standard Proof.}~
TODO: supply the compact proof.
\end{proof}

\begin{proof}
\textbf{Detailed Learning Proof.}~
TODO: supply the detailed learning proof.
\end{proof}

\begin{remark*}[Proof structure]
Assume the two conditionals are true. If \(\varphi\) is true, then \(\psi\) is
true, and hence \(\chi\) is true.
\end{remark*}

\begin{dependencies}
\begin{itemize}
  \item \hyperref[def:valid-argument-form-propositional-logic]{Valid Argument Form}
  \item \hyperref[def:logical-consequence-propositional-logic]{Logical Consequence}
\end{itemize}
\end{dependencies}

\clearpage
% END GENERATED ARTIFACT: prf:hypothetical-syllogism-validity-propositional-logic
